\subsection{E(3)-equivariant message passing, readout heads, and architecture variants}
\label{subsec:architecture}

The network is composed of node encoders, edge encoders, a stack of equivariant
message-passing blocks, and matrix-specific readout heads. Node encoders embed
chemical species into scalar channels. Edge encoders combine edge-type
embeddings, radial distance embeddings, and spherical harmonics into equivariant
edge features. If $\bm{h}_i^{(t)}$ denotes node features and
$\bm{e}_{ij}^{(t)}$ denotes edge features at layer $t$, a generic Mandala
message-passing step may be written as
\begin{align}
\bm{m}_{ij}^{(t)} &=
\Phi_{\mathrm{edge}}^{(t)}
\left(
\bm{h}_i^{(t)}, \bm{h}_j^{(t)}, \bm{e}_{ij}^{(t)},
Y(\hat{\bm{r}}_{ij}), \varphi(r_{ij})
\right), \\
\bm{h}_i^{(t+1)} &=
\Phi_{\mathrm{node}}^{(t)}
\left(
\bm{h}_i^{(t)},
\mathrm{Agg}_{j \in \mathcal{N}(i)} \, \bm{m}_{ij}^{(t)}
\right),
\end{align}
where $\Phi_{\mathrm{edge}}^{(t)}$ and $\Phi_{\mathrm{node}}^{(t)}$ are
equivariant maps assembled from tensor products, equivariant linear layers,
radial MLPs, parity-aware nonlinearities, and optional residual connections,
while $\mathrm{Agg}$ denotes a configurable aggregation rule such as summation
or attention.

The readout stage maps final latent features to the irrep content required by
each matrix target. For a target operator $X$, the head predicts irrep vectors
\begin{equation}
\widehat{\bm{x}}_{ij}^{\bm{L},X}
=
\Psi_X\!\left(
\bm{e}_{ij}^{(L)},
\bm{h}_i^{(L)},
\bm{h}_j^{(L)},
\varphi(r_{ij})
\right),
\end{equation}
followed by inverse irrep-to-block conversion
\begin{equation}
\widehat{X}_{ij}^{\bm{L}}
=
Q_{a_i a_j}^{\top}\widehat{\bm{x}}_{ij}^{\bm{L},X}.
\end{equation}
Because heads share the same equivariant latent representation, Mandala can
predict Hamiltonian, overlap, and density matrices jointly with only modest
additional readout cost.

An important feature of the framework is that the architecture is not fixed to
one canonical model. Mandala exposes a broad family of variants within the same
software interface. Search dimensions include the hidden irrep content, the
angular cutoff $\ell_{\max}$, the number of message-passing layers, the radial
basis size, the depth of the neck and head MLPs, and the choice of tensor
product implementation. The message-passing blocks support alternative edge
updates (tensor-product, concatenation-based, and replacement-style variants),
alternative node updates (sum, concatenation, and replacement variants),
alternative aggregation mechanisms (including attention-style aggregation),
pre- and post-linear transformations, optional self-connections, and optional
node and edge residual paths.

The encoder and head stages are likewise configurable. Edge encoders support
various variants, together with optional tensor-square
feature expansion of spherical-harmonics. Readout heads support direct
equivariant projections, deeper equivariant MLP heads, optional tensor-square
preconditioning of hidden features, scalar-specific auxiliary branches, optional
log-scale or magnitude factorization for output amplitudes, optional use of node
embeddings for self-edges, and separate treatment of diagonal, shifted-self,
and off-diagonal blocks. Because all of these variants share the same
block-sparse data model and irrep mapping layer, architecture search can be
performed without changing the surrounding training and evaluation workflow.
